\section{Память и соблюдение лимита}
\label{sec:memory}

\subsection{Принцип out-of-core}
Данные делятся по размеру. Структуры $\bigO(n)$ (по числу вершин) живут в ОЗУ. Структуры $\bigO(m)$
(по числу рёбер) живут на диске и читаются потоком. Поскольку <<гигабайт>> --- это именно рёбра,
потребление памяти определяется числом вершин, а не рёбер. Отсюда следствие, заметное на практике:
размер RSS зависит от числа вершин, а не от размера CSV-файла. Меньший по объёму файл с большим
числом вершин требует больше памяти, чем больший файл с малым числом вершин.

\subsection{Из чего складывается память}
\label{sec:buckets}
Потребление ОЗУ складывается из трёх частей.
\begin{itemize}[itemsep=3pt]
  \item \textbf{Корзина A --- общие векторы $\bigO(n)$.} Три \texttt{double[n]} ($s_{\mathrm{old}}$,
        $s_{\mathrm{new}}$, $\mathrm{contrib}$) и несколько \texttt{int[n]} (\texttt{sourcesPtr},
        \texttt{outdeg}, \texttt{originalIds}) --- это примерно $36$ байт на вершину. Для
        $n\approx 4{,}8$ млн --- около $170$~МБ. Эта часть общая для всех потоков и от их числа не
        зависит.
  \item \textbf{Корзина B --- попоточные буферы.} Один буфер чтения на поток ($64$~КиБ). Всего
        $P\cdot B$ плюс стеки потоков --- единицы МБ. Это единственная часть, растущая с числом
        потоков.
  \item \textbf{Корзина C --- накладные расходы JVM.} Metaspace, кэш кода, структуры сборщика ---
        около $20$--$30$~МБ. Serial GC держит эту часть малой.
\end{itemize}
Основной объём --- корзина A --- общий и не уменьшается с ростом числа потоков. С потоками растёт
только небольшая корзина B.

\subsection{Лимит памяти и его проверка}
\label{sec:cap}
Ограничение памяти задаётся снаружи, а программа подстраивает под него своё поведение. Доступный объём
она узнаёт из максимального размера кучи (значение \texttt{-Xmx}, оно же видно как лимит cgroup в
контейнере) и под него подбирает размер корзин $M$ при предобработке. Один рычаг (\texttt{-Xmx} или
cgroup) задаёт и реальный потолок памяти, и параметры предобработки.

Чтобы симулировать ограничение из условия и проверить, что оно действительно соблюдается, использованы
два механизма.
\begin{itemize}[itemsep=3pt]
  \item \textbf{Лимит кучи JVM (\texttt{-Xmx}).} Ограничивает кучу Java. Нужен, но его одного мало:
        вся память процесса --- это куча плюс стеки, metaspace, кэш кода, буферы вне кучи и страничный
        кэш. Часть из этого лежит вне \texttt{-Xmx}.
  \item \textbf{Лимит всего процесса через cgroup~v2 (основной).} JVM запускается внутри ограничения
        cgroup, и ядро ОС ограничивает всё сразу, снимая процесс по OOM при нарушении:
        \begin{verbatim}
    systemd-run --user --scope -p MemoryMax=256M -p MemorySwapMax=0 \
        java -Xmx128m -XX:+UseSerialGC ...
        \end{verbatim}
        Своп отключён (\texttt{MemorySwapMax=0}), иначе куча незаметно вытеснялась бы в своп и лимит
        стал бы фикцией.
\end{itemize}
Само завершение запуска под лимитом cgroup и есть доказательство: при превышении процесс был бы снят
ядром. Дополнительно в конце запуска печатается пиковый RSS (значение \texttt{VmHWM} из
\texttt{/proc/self/status}) --- получается конкретное число вида <<пик RSS\,$=168$~МиБ при лимите
$256$~МиБ>>.

\paragraph{Деталь про \texttt{-Xmx}.} Он ограничивает только кучу. Поверх неё JVM тратит ещё
$60$--$80$~МБ на нативную часть (сам \texttt{libjvm}, кэш кода, стеки). Поэтому, чтобы удержать RSS
под лимитом $C$, \texttt{-Xmx} стоит брать примерно на $80$~МБ меньше $C$, а не равным $C$.

\subsection{Контрольный контрпример}
Чтобы показать ценность out-of-core, в проект включён намеренно наивный вариант, который держит весь
граф в памяти (вложенные списки в ОЗУ). Под тем же лимитом cgroup он падает, а потоковый движок
завершается. Падает он двумя способами. При куче, равной потоковой, ему не хватает кучи Java
(\texttt{OutOfMemoryError}). При щедрой куче (\texttt{-Xmx1g}) его снимает ядро по OOM (сигнал KILL,
код $137$). Алгоритм и лимит те же, отличается только способ хранения.

\subsection{Ограничения подхода}
\label{sec:limits}
\begin{itemize}[itemsep=3pt]
  \item В ОЗУ обязаны помещаться структуры $\bigO(n)$ --- векторы значений и указатели
        (раздел~\ref{sec:scope}). Если не влезают даже они, это уже другая задача (выгрузка самих
        векторов на диск, как в GraphChi PSW), и она вне рамок работы.
  \item На очень больших графах узкое место --- скорость последовательного чтения с диска (файл
        \texttt{sources} перечитывается на каждой итерации). Компактный формат поднимает этот потолок,
        но саму природу ограничения не убирает.
  \item $128$~МБ --- это тесно для JVM. Нативный пол ($\approx 70$~МБ) плюс векторы $\bigO(n)$ при $n$
        около $2{,}7$ млн ($\approx 108$~МБ) уже превышают $128$~МБ. При таком жёстком лимите это
        аргумент в пользу C++. На Java тот же подход работает при чуть большем лимите.
\end{itemize}
